% ---------------------------------------------------------------------------
% Related work, compressed to the four traditions and what each leaves open.
% The economics paragraph is load-bearing here: the asymmetry this paper
% measures is the testable form of the switching-cost claim.
% ---------------------------------------------------------------------------

\section{Background and Related Work}
\label{s:sec:background}

Four traditions bear on the question of \S\ref{s:sec:intro} and none of them answers it.
Formal-language theory supplies the frame; compiler correctness supplies the standard of
proof; the expressiveness literature supplies the notion of ``cannot say''; and the
economics of switching costs supplies both the reason anybody should care and, as it turns
out, the shape of the answer.

\paragraph{Query languages as formal languages.}
Rich's treatment of query languages as formal languages is the licence for the approach: a
query language is a formal language, so grammars, automata and decidability apply to it,
and which class a language falls in decides what can be settled
mechanically~\citep{rich2008}. §48.1 is where the book says so directly, in a short
applications note pointing at published SQL grammars; the
regular/context-free/decidable classification this paper leans on is the book's frame
rather than that section's. An unambiguous grammar gives each query exactly one parse
tree, which is the precondition for a translation being well defined; an LALR(1) one makes
the parse linear and the parser a table, so a corpus can be processed in bulk rather than
sampled~\citep{aho2006}. The translation itself has a classical form: a syntax-directed
transduction associates an output with each production and composes them over the parse
tree~\citep{lewis1968}, and where a node's output depends on information held elsewhere,
Knuth's attribute grammars are the standard account~\citep{knuth1968} --- which this
problem needs, because two of the attributes a translation requires are not in the Datadog
query text at all. Regular languages then enter directly, and it is where the sharpest
negative results live: both languages filter tag values with a pattern language, a glob in
one and a regular expression in the other, so whether one pattern has a counterpart is a
question about two classes of regular languages (\S\ref{s:sec:regular}).

\paragraph{The two languages.}
PromQL is documented and its reference parser is open source, so a grammar can be checked
against an implementation rather than inferred~\citep{promql}. It is a vector-valued
expression language: selectors denote sets of series, range selectors denote windows,
functions reduce windows, and aggregation operators reduce across series. Counters are
stored as running totals and reset handling lives inside the range functions. Datadog's
metric query language is not documented as a grammar; what is published is a reference
manual with examples, enough to infer the syntax and not to verify it. The important
difference is not syntactic: Datadog stores a counter as per-interval deltas, and a query
is evaluated first by time and then by space~\citep{datadog,ddnested} --- a rollup reduces
each series alone and a space aggregator combines the results. Neither the rollup interval
nor the metric's type is part of the query text, which is the subject of the complete
paper's largest finding and of \S\ref{s:sec:asymmetry} here.

\paragraph{What a translation guarantees.}
SQL dialect translation is the closest established practice. Apache Calcite parses SQL
into a relational algebra, optimises, and generates dialect-specific SQL on the way
out~\citep{begoli2018}; SQLGlot does the same across more than twenty
dialects~\citep{sqlglot}. Both are engineering successes and neither offers a statement of
what a translation preserves: their criterion is that the output parses and behaves as
users expect, and where a construct has no counterpart the usual outcome is a best-effort
rewrite that leaves no record of the approximation. That is a reasonable engineering choice
and a poor measurement instrument, which is why \S\ref{s:sec:framework} requires a
translator to be total with every failure a named $\bot$.

Converters for this language pair exist, and we have found no published result among them
--- a negative about a literature, which is the weakest kind of claim this paper makes and
is stated as a search rather than as a fact. Chronosphere describes the closest thing to
this paper's method, a Datadog grammar written from scratch with queries parsed to an AST
and emitted through the open-source Prometheus grammar, and reports getting ``about
$90\%$'' of the way~\citep{marcopoli2024}; what the percentage is a percentage of, and what
a converted query is meant to preserve, are not stated. OpenObserve's dashboard migrator
marks each object migrated, needing review, or skipped~\citep{openobserve} --- a coarse
form of the typed failure \S\ref{s:sec:framework} requires, reported without counts. A
third converter is closer to this paper than either: a public repository with no
accompanying publication, which translates Datadog dashboards and monitors into PromQL and
reports reaching $92.4\%$ of $14{,}172$ distinct queries and monitors~\citep{dd2dash0}.
Those figures are its own and not this paper's, and they are read from a cached copy of its
README because the repository is a feature branch. It arrives independently at three of
this paper's devices: each of its rules carries one of four claims, each translated item
carries a ledger of the assumptions its claim rests on, and anything without a rule is
refused by the name of the construct rather than approximated. That convergence is the
useful part of finding it --- these are not this paper's inventions but what the problem
forces on anyone unwilling to round a translation up. What none of them publishes is a
grammar for the Datadog side, or a measurement of what is impossible as against unbuilt,
which is what leaves this paper's claims standing.

\paragraph{The standard of proof, and the form of a negative.}
Compiler verification supplies the first. McCarthy and Painter proved a compiler for
arithmetic expressions correct by relating source and target denotations~\citep{mccarthy1967},
and CompCert carries the idea through a realistic C compiler with machine-checked
proofs~\citep{leroy2009}. That is the shape of this paper's soundness theorems though not
their rigour --- the proofs here are on paper and for a fragment --- and it differs in one
way that is not incidental: a verified compiler may assume source and target run on the
same input, and here they do not. Database theory supplies the adjacent decision problem,
equivalence and containment of conjunctive queries being decidable by the homomorphism
characterisation~\citep{chandra1977,abiteboul1995}; it does not transfer, since the
quantities here are numeric aggregates over time windows and exactness has to be stated up
to floating point. The borrowing is the habit: state the equivalence you mean, then decide
it.

``PromQL cannot say this'' is also a claim with a literature behind it. Felleisen
formalises expressive power by asking whether a construct can be eliminated by a local,
meaning-preserving rewriting~\citep{felleisen1991}, and Mitchell relates expressiveness to
observational equivalence~\citep{mitchell1993}. What this paper borrows is the
\emph{distinction} rather than the criterion: \S\ref{s:sec:regular}'s results are of a
blunter kind, containment facts about two classes of regular languages where the target has
no term with the denotation at all. They are the only claims here that hold regardless of
how much engineering is applied.

\paragraph{Observability, standardisation, and the asymmetry.}
The closest observability work compares implementations of \emph{one} language rather than
translating between two. Volz's compliance testing found that vendors advertising PromQL
support disagree with Prometheus on a generated conformance suite~\citep{volz2020} --- this
paper's problem in miniature, and a warning: if one language's implementations already
differ, a translation between two needs its soundness stated rather than assumed. The
practice of verifying a candidate against a real server rather than a model is taken from
that result. Standardisation approaches it from the other end: OpenTelemetry has largely
settled the \emph{data} question~\citep{otel}, leaving queries as the remaining lock-in
surface, and the CNCF's Observability TAG has an active effort towards a common query
language~\citep{sharma2025} --- which needs exactly the evidence this paper produces, namely
which constructs are genuinely hard to translate and which merely have not been.

Switching costs themselves are well developed in industrial organisation. Farrell and
Klemperer survey lock-in and its competitive effects~\citep{farrell2007}, and Shapiro and
Varian give the lock-in cycle its standard informal treatment for information
goods~\citep{shapiro1999}. The structural point matters more here than any quantity, and it
is what this paper ends up measuring: what retains a customer is not that leaving is
impossible but that \emph{the terms of return differ from the terms of entry}, so a
decision made cheaply cannot be unmade cheaply. That is an asymmetry claim, it is testable,
and \S\ref{s:sec:asymmetry} is the test. The concrete case that motivates the question is
on the public record~\citep{orosz2023}.

\paragraph{What this paper adds.}
Each tradition leaves the same thing open. The transpilers translate without stating what
is preserved, and the migration tools for this pair report neither what they preserve nor
what they refuse; the verification work states what is preserved but assumes a shared
input; the expressiveness literature gives the form of a negative result but is rarely
pointed at query languages; the observability work compares implementations of one
language, and the standardisation work needs evidence about which differences are
essential; the economics predicts an asymmetry without measuring one. The contribution is
the combination: soundness stated relative to an ingestion mapping, so that two
differently-stored systems can be compared at all; a translator whose refusals are typed
and counted rather than papered over; non-expressibility proved where it holds; and all of
it measured on a corpus of queries nobody wrote for the purpose.
